\documentclass[12pt,a4paper,oneside]{article}
\usepackage[brazil]{babel}
\usepackage[utf8]{inputenc}
\usepackage[dvipsnames]{xcolor}
\usepackage{listings}
\usepackage{wrapfig}
\usepackage{olymp}

\setcounter{problem}{2}

\begin{document}

\begin{problem}{Notas de prova}{notas}{2}

Certo professor terminou de corrigir todas as provas de sua turma e está com as notas anotadas numa ordem qualquer, como mostrado no exemplo abaixo, à esquerda. Como as notas são lançadas no Sapiens em ordem de matrícula, facilitaria muito seu trabalho se a lista de notas estivesse ordenada pelo número de matrícula, como a lista da direita!\\

\hspace{3cm}\begin{minipage}{.3\linewidth}
\begin{verbatim}
Matrícula Nota
  60001    20
  61234    23
  59999    18
  56789    19
  66666    22
\end{verbatim}
\end{minipage}\hspace{3cm}
\begin{minipage}{.3\linewidth}
\begin{verbatim}
Matrícula Nota
  56789    19
  59999    18
  60001    20
  61234    23
  66666    22
\end{verbatim}
\end{minipage}\\

Ajude o professor fazendo um programa que ordena as notas pela matrícula.


% \Notes

% Uma seleção ganha 3 pontos para cada vitória e 1 ponto para cada empate. Derrotas não alteram sua pontuação.

\InputFile

A entrada começa com uma linha contendo um número $T$, que é o número de alunos da turma. Em seguida vem $T$ linhas, cada uma com dois inteiros $M_i$ e $N_i$, que são o número de matrícula e a nota do $i$-ésimo aluno da lista, respectivamente. Restrições: $1 \leq T \leq 100$; $1 \leq M_i \leq 99999$; $0 \leq N_i \leq 30$. Não há número de matrícula repetido.

\OutputFile

Seu programa deve gerar $T$ linhas de saída, cada uma contendo o número de matrícula e a nota de um dos alunos. A saída deve estar ordenada pelo número de matrícula.

% \Notes
% Preste atenção aos tipos das variáveis, pois $F(90) \approx 2^{61}$.

\Example

\begin{example}
\exmp{
5
60001 20
61234 23
59999 18
56789 19
66666 22
}{
56789 19
59999 18
60001 20
61234 23
66666 22
}%
\end{example}

\begin{example}
\exmp{
4
20000 18
19000 15
18000 20
17000 15
}{
17000 15
18000 20
19000 15
20000 18
}%
\end{example}


\begin{example}
\exmp{
2
54321 9
12345 20
}{
12345 20
54321 9
}%
\end{example}

\textbf{Dica}: crie uma \texttt{struct Aluno} com os campos matrícula e nota; armazene os dados de entrada em um vetor de \texttt{Aluno} e então ordene o vetor de acordo
com as matrículas dos alunos.

%Comentários sobre o exemplo, se necessário.

\end{problem}

\end{document}
